% Standalone problem document, generated from the adjacent README.md.
% Compile with XeLaTeX. Mathematical content is maintained in Markdown.
\documentclass[11pt,a4paper]{article}
\usepackage[fontsize=10.5pt]{scrextend}
\usepackage[margin=22mm,headheight=15pt,headsep=9mm,footskip=11mm]{geometry}
\usepackage{fontspec}
\setmainfont{texgyrepagella}[Extension=.otf,UprightFont=*-regular,BoldFont=*-bold,ItalicFont=*-italic,BoldItalicFont=*-bolditalic]
\setsansfont{texgyreheros}[Extension=.otf,UprightFont=*-regular,BoldFont=*-bold,ItalicFont=*-italic,BoldItalicFont=*-bolditalic]
\setmonofont{texgyrecursor}[Extension=.otf,UprightFont=*-regular,BoldFont=*-bold,ItalicFont=*-italic,BoldItalicFont=*-bolditalic,Scale=MatchLowercase]
\usepackage{amsmath,amssymb}
\usepackage{unicode-math}
\setmathfont{texgyrepagella-math.otf}
\usepackage{xcolor,microtype,enumitem,needspace,fancyhdr}
\usepackage{bookmark}
\definecolor{ink}{HTML}{17344B}
\definecolor{accent}{HTML}{197D80}
\definecolor{muted}{HTML}{596773}
\hypersetup{colorlinks=true,linkcolor=accent,urlcolor=accent,pdftitle={TR-05 - Infinite
mean condition number in every identifiable tensor
format},pdfauthor={Open Problems in Numerical Linear Algebra}}
\urlstyle{same}
\setlength{\parindent}{0pt}
\setlength{\parskip}{5pt plus 1pt minus 1pt}
\linespread{1.04}
\setlength{\emergencystretch}{3em}
\widowpenalty=10000
\clubpenalty=10000
\displaywidowpenalty=10000
\allowdisplaybreaks[1]
\setlist{nosep,leftmargin=1.5em}
\providecommand{\tightlist}{\setlength{\itemsep}{0pt}\setlength{\parskip}{0pt}}
\providecommand{\pandocbounded}[1]{#1}
\pagestyle{fancy}
\fancyhf{}
\fancyhead[L]{\sffamily\footnotesize\color{muted}OPEN PROBLEMS IN NUMERICAL LINEAR ALGEBRA}
\fancyhead[R]{\sffamily\bfseries\footnotesize\color{accent}TR-05}
\fancyfoot[L]{\parbox[b]{0.9\textwidth}{\sffamily\fontsize{8}{10}\selectfont\color{muted}Editorial ratings; a bounded search cannot certify that a problem remains unsolved.\\Literature check: 2026-09-10}}
\fancyfoot[R]{\sffamily\footnotesize\color{muted}\thepage}
\renewcommand{\headrulewidth}{0.3pt}
\renewcommand{\headrule}{\hbox to\headwidth{\color{accent}\leaders\hrule height \headrulewidth\hfill}}
\makeatletter
\renewcommand{\section}{\Needspace{5\baselineskip}\@startsection{section}{1}{0pt}{12pt plus 2pt minus 2pt}{4pt}{\sffamily\bfseries\large\color{ink}}}
\renewcommand{\subsection}{\Needspace{4\baselineskip}\@startsection{subsection}{2}{0pt}{9pt plus 1pt minus 1pt}{3pt}{\sffamily\bfseries\normalsize\color{ink}}}
\makeatother
\setcounter{secnumdepth}{0}
\begin{document}
{\sffamily\small\color{muted}Tensor computations\par}
\vspace{3pt}
{\raggedright\hyphenpenalty=10000\exhyphenpenalty=10000\sffamily\bfseries\fontsize{21}{24}\selectfont\color{ink}Infinite
mean condition number in every identifiable tensor format\par}
\vspace{7pt}
{\sffamily\small\textbf{Difficulty:} challenging\quad\textbf{Importance:} interesting
to the community\par}
{\sffamily\small\bfseries\color{accent}Status: Partially resolved\par}
\vspace{4pt}
{\color{accent}\hrule height 0.8pt}
\vspace{6pt}
\textbf{Rating rationale:} Challenging because extending the
mean-divergence theorem requires controlling singular decomposition
geometry beyond the existing identifiability hypothesis; community
importance comes from average numerical conditioning of CP
decomposition.

\subsection{Context and notation}\label{context-and-notation}

Fix \(d\ge3\), \(n_j\ge2\), and \(r\ge2\). Assume generic complex
identifiability: outside a proper algebraic exceptional set, a complex
tensor of rank \(r\) in this format has a unique unordered collection of
\(r\) rank-one summands. Let \(M_r\) be the smooth identifiable locus of
real rank-\(r\) tensors, with induced Euclidean volume \(dV\). The
random input has density

\[
d\mu(A)=Z^{-1}e^{-\|A\|_F^2/2}\,dV(A).
\]

For the addition map \(\Phi(a_1,\ldots,a_r)=\sum_i a_i\) on rank-one
tensors, let \(\Psi\) be a local inverse at \(A\). Use Frobenius norms
and their product norm. Define

\[
\kappa(A)=\|D\Psi(A)\|_2,\qquad
\kappa_{\rm ang}(A)=\|D(p^{\times r}\circ\Psi)(A)\|_2,
\quad p(a)=a/\|a\|_F.
\]

Values on measure-zero exceptional sets do not affect the expectations.
This samples tensors by volume, not independent Gaussian summands.

\subsection{Problem statement}\label{problem-statement}

Under this probability model, prove or disprove

\[
\mathbb E_\mu\kappa(A)=\infty
\]

for every admissible format and \(r\ge2\). Rank two is already proved.
The unresolved task removes the additional smaller-format
identifiability assumption used for the higher-rank theorem; it must not
be counted again for parameter ranges covered by that theorem.

\subsection{Reference}\label{reference}

Carlos Beltrán, Paul Breiding, and Nick Vannieuwenhoven,
\href{https://doi.org/10.1007/s10208-022-09551-1}{\emph{The Average
Condition Number of Most Tensor Rank Decomposition Problems Is
Infinite}}, \emph{Foundations of Computational Mathematics} 23 (2023),
433--491: Definition 1, Assumption 1, equation (5), Theorems 1--2, and
Conjecture 1. The \href{https://arxiv.org/pdf/1903.05527}{author
preprint} labels the conjecture 1.8.

\subsection{Status check --- 2026-09-10}\label{status-check-2026-09-10}

Rechecked the \href{https://doi.org/10.1007/s10208-022-09551-1}{journal
version, Theorems 1--2 and Conjecture 1}, and searched for subsequent
average-condition-number resolutions. Rank two and the higher-rank
formats satisfying the theorem's extra smaller-format identifiability
assumption are proved portions of the displayed family; removing that
assumption remains unresolved. No later general proof was located. The
latest explicit conjecture checked remains the 2023 journal version, so
the later-status evidence is a bounded search.
\end{document}
